\chapter{Solution}

The implemented solution is a file-based research pipeline inside the qyl repository. The corpus configuration is stored in \texttt{research/arqio/corpus/targets.json}. Result schemas are stored under \texttt{research/arqio/schemas}. The main corpus runner is \texttt{eng/automation/arqio/run-corpus.sh}, which delegates to a Python implementation for structured JSON output. Coverage is handled by \texttt{eng/automation/coverage/run-dotnet-coverage.sh}. Semantic-convention verification is handled by \texttt{eng/automation/semconv/verify-qyl-semconv.sh}. Thesis quality checks are handled by \texttt{eng/automation/thesis/check-thesis.sh}.

The pipeline preserves raw evidence instead of summarizing away failures. Each target receives an output directory under \texttt{research/arqio/results/<run-id>/<target-id>/iteration-<n>}. Logs include restore, build, test, coverage, classification, and Docker compose configuration checks when telemetry is requested. The aggregate result is written to \texttt{aggregate-result.json}. Tables are generated from that aggregate rather than hand-written, which reduces the risk of thesis text drifting away from measured data.

\begin{figure}[h]
\centering
\begin{verbatim}
Target repo -> restore/build/test -> coverage collector
            -> runtime telemetry check -> Arqio classifier
            -> aggregate result JSON -> Rego rubric input
            -> thesis tables -> thesis.pdf
\end{verbatim}
\caption{Local evaluation pipeline.}
\end{figure}

The OpenTelemetry path is conservative. The harness does not invent semantic attributes and does not publish packages. It verifies local semantic-convention projects and records whether Docker-based telemetry configuration is available. If runtime spans are not collected, the result remains unavailable. This is a deliberate no-fabrication rule: telemetry absence weakens the thesis, but fake spans would invalidate it.
